% Budget : 10 pages — voir rapport/PLAN.md
\chapter{Cadre général du projet}

\section{Introduction}

Ce chapitre pose le cadre dans lequel s'inscrit ce projet de fin d'études. Nous présentons d'abord
l'entreprise d'accueil, puis le contexte, la problématique, les objectifs et le périmètre du projet.
Nous analysons ensuite les solutions existantes afin de situer notre approche par rapport à elles.
Nous terminons par la méthodologie de développement retenue.

\section{Présentation de l'organisme d'accueil}

Ce projet a été réalisé au sein de l'entreprise \textbf{ASM (All Soft Multimedia)}, une société de
services numériques tunisienne fondée en 2007. Éditrice et intégratrice de solutions logicielles,
ASM conçoit des sites web, des applications mobiles, des systèmes de point de vente ainsi qu'un
progiciel de gestion intégré propriétaire, \textbf{DUX}, distribué auprès de plus de 4500 clients.
Son siège se trouve à Sfax, avec une présence à Sousse et Tunis ainsi que des filiales à Alger et à
Dakar~\cite{asm-apropos}.

\begin{figure}[H]
  \centering
  \includegraphics[width=0.35\textwidth]{logo-asm}
  \caption{Logo de l'entreprise ASM (All Soft Multimedia)~\cite{asm-apropos}}
\end{figure}

Ce positionnement est directement pertinent pour ce projet : ASM n'est pas seulement le
commanditaire, elle est elle-même éditrice d'un ERP. Une plateforme de suivi de livraison qui
s'intègre à un ERP sans s'y enfermer profite à ses clients quel que soit le progiciel qu'ils
exploitent, et son architecture, bâtie sur un contrat d'adaptateur commun à tout ERP, est conçue pour qu'un autre système
puisse s'y brancher sans remettre en cause le reste de la plateforme.

\begin{table}[H]
\centering
\caption{Fiche d'identité d'ASM}
% Pas de traits verticaux : convention observée dans le mémoire de référence, où tous les
% tableaux comparatifs ne sont séparés que par des filets horizontaux.
\begin{tabularx}{\textwidth}{L{4cm}X}
\hline
\rowcolor{isimsblue!10}
\textbf{Nom de l'entreprise} & ASM (All Soft Multimedia) \\
\hline
\textbf{Secteur d'activité} & Édition et intégration de logiciels : ERP (DUX), sites web,
applications mobiles, systèmes de point de vente \\
\hline
\textbf{Date de création} & 2007 \\
\hline
\textbf{Siège social} & Sfax, Tunisie \\
\hline
\textbf{Présence} & Sfax, Sousse, Tunis (Tunisie) ; Alger (Algérie) ; Dakar (Sénégal) \\
\hline
\textbf{Site web} & \url{https://asm-tunisie.com} \\
\hline
\end{tabularx}
\end{table}

\section{Présentation du projet}

\subsection{Contexte du projet}

Un ERP, quel qu'il soit, gère efficacement les commandes, les stocks et la facturation, mais
s'arrête à la porte de l'entrepôt. Tout ce qui se passe ensuite, l'affectation d'un livreur, le suivi
de sa tournée, la collecte de la preuve de remise et un éventuel retour, reste hors de sa portée, ou
suppose de développer un module spécifique à chaque ERP concerné. Or les entreprises que
sert ASM n'exploitent pas toutes le même système, et ce choix ne devrait pas
déterminer si elles peuvent, ou non, suivre leurs livraisons.

\subsection{Problématique}

La problématique de ce projet peut être résumée ainsi :

\begin{quote}
\bfseries\itshape
Comment concevoir une plateforme de suivi de livraison exploitable par plusieurs entreprises
clientes, indépendante de l'ERP utilisé par chacune, et capable de fonctionner sur le terrain même
en l'absence de connexion réseau~?
\end{quote}

\subsection{Objectifs du projet}

L'objectif principal est de couvrir l'ensemble du cycle de vie d'une livraison, depuis l'import de
la commande jusqu'à la preuve de sa remise, à travers :

\begin{itemize}
  \item une architecture \textbf{multi-tenant}, où chaque entreprise cliente dispose de ses propres
  données isolées au sein d'une même plateforme.
  \item une intégration ERP \textbf{agnostique}, opérationnelle sur Odoo et ERPNext par un
  adaptateur commun plutôt qu'un développement spécifique à chacun, et conçue pour qu'un autre ERP
  puisse s'y brancher sans toucher au reste du système.
  \item une application mobile chauffeur fonctionnant en \textbf{mode hors ligne}, pour que la
  perte de réseau sur le terrain n'interrompe pas la livraison.
  \item un suivi en \textbf{temps réel} des tournées, de la planification jusqu'à la remise au
  client final.
  \item une \textbf{chaîne d'encaissement} traitée comme un circuit de responsabilité, de la
  collecte des espèces à la porte du client jusqu'à leur remise au dépôt et l'arbitrage des
  écarts.
  \item un \textbf{assistant interne} auquel un opérateur pose ses questions en langage naturel,
  et qui répond sur l'activité de son entreprise sans jamais lui donner accès à ce que son rôle
  ne l'autorise pas à consulter.
\end{itemize}

\subsection{Périmètre du projet}

Le projet couvre le cycle complet d'une livraison, de l'import à la planification de tournée, puis
l'exécution terrain, la preuve de livraison, l'encaissement et les retours, pour un nombre quelconque d'entreprises
clientes. Il se limite au dernier kilomètre : la gestion des stocks, la facturation et la production
restent du ressort de l'ERP, que la plateforme consulte et met à jour sans s'y substituer.

\section{Étude et critique de l'existant}

\subsection{Aperçu des outils existants}

Avant de concevoir la solution, nous avons examiné les outils déjà utilisés pour le suivi de
livraison du dernier kilomètre.

\begin{figure}[H]
  \centering
  \includegraphics[width=0.15\textwidth]{logo-onfleet}
  \caption{Logo d'Onfleet~\cite{onfleet-site}}
\end{figure}

\textbf{Onfleet}~\cite{onfleet-site} est une plateforme SaaS de gestion de tournées et de suivi de
chauffeurs, reconnue pour son optimisation d'itinéraires et son suivi en temps réel. Elle ne
s'intègre toutefois à un ERP que par des connecteurs génériques ou un développement sur mesure. Elle propose un
mode hors ligne, mais désactivé par défaut : sans lui, un chauffeur ne dispose que de deux minutes
pour démarrer ou terminer une tâche avant d'être bloqué faute de réseau~\cite{onfleet-connectivity}.
Sa documentation publique ne décrit pas de mécanisme d'encaissement en espèces ni de remise de
caisse.

\begin{figure}[H]
  \centering
  \includegraphics[width=0.15\textwidth]{logo-bringg}
  \caption{Logo de Bringg~\cite{bringg-site}}
\end{figure}

\textbf{Bringg}~\cite{bringg-site} est une plateforme de logistique de livraison orientée grands
comptes, proposant orchestration de flotte et visibilité client. Son application chauffeur reste
utilisable hors ligne pour les tâches déjà assignées~\cite{bringg-connectivity}, et son API prévoit
la prise d'un paiement par le livreur ainsi qu'un montant restant dû sur la
commande~\cite{bringg-driver-actions}. Sa richesse fonctionnelle a pour contrepartie une mise en
œuvre lourde, pensée pour un déploiement par ERP ou canal de vente plutôt que pour une intégration
légère et multi-ERP.

\textbf{Le module de livraison natif d'un ERP} (par exemple Odoo Inventory) constitue la troisième
approche possible. Il a l'avantage d'être déjà connecté aux données de commande et de stock, mais il
enferme l'entreprise dans cet ERP précis : le changer, ou en exploiter plusieurs à la fois, ce qui
est un cas concret chez les clients d'ASM, n'est pas pris en charge.

\subsection{Synthèse comparative et solution proposée}

Le fonctionnement hors connexion, longtemps distinctif, est aujourd'hui proposé par les plateformes
matures du marché. Ce qui les sépare de notre approche tient à deux points : la profondeur de
l'intégration ERP, générique chez elles alors que le projet vise un contrat commun réellement
implémenté sur plusieurs ERP, et le traitement de l'encaissement à la livraison. Sur ce second
point, la distinction porte moins sur la collecte du paiement, que Bringg documente à travers une
action chauffeur de prise de paiement et un montant restant dû, que sur ce qui vient après : la remise de
la caisse du livreur, son comptage contradictoire et l'arbitrage des écarts. Ce maillon est
déterminant sur le marché tunisien, où le
paiement en espèces à la remise reste le mode dominant.

Un troisième écart tient à la nature de l'assistance intégrée. Onfleet a doté son tableau de bord
d'un assistant conversationnel en avril~2026, inclus dans tous ses abonnements~; sa documentation le
présente comme une aide à l'usage du produit~: il répond aux questions de prise en main et de
dépannage~\cite{onfleet-assistant}. Bringg, de son côté, ne documente aucun assistant
conversationnel et concentre son intelligence artificielle sur la prédiction et l'optimisation des
tournées~\cite{bringg-ptos}. Aucun des deux ne documente donc un assistant capable d'interroger les
données d'exploitation du client en langage naturel.

C'est précisément la fonction retenue ici, avec deux garde-fous qui en font autre chose qu'un agent
conversationnel branché sur une base. D'abord, l'assistant \textbf{hérite des droits de celui qui
l'interroge}~: il ne dispose d'aucun compte propre, et deux utilisateurs de rôles différents
n'obtiennent pas la même réponse à la même question. Ensuite, il \textbf{ne peut que lire}~: aucune
de ses opérations ne modifie l'état du système. Il cite enfin ses sources, de sorte que chaque
réponse reste vérifiable par celui qui la reçoit.

\begin{table}[H]
\centering
\caption{Comparatif des solutions existantes, d'après leur documentation publique}
\label{tab:comparatif-existant}
% Ce critère a d'abord porté sur la multi-tenance, puis sur l'auto-hébergement : les deux étaient
% erronés. Onfleet et Bringg sont des SaaS multi-clients par construction, et la solution proposée
% est elle aussi hébergée par son éditeur (ASM) plutôt que par le client.
%
% La colonne COD a elle-même dû être corrigée : elle portait « Non » pour Bringg, alors que sa
% documentation API expose bien un montant restant dû et une action chauffeur de prise de paiement.
% D'où la formulation « Partiel » et le libellé « Non documenté » plutôt que « Non » — une absence
% dans la documentation publique ne prouve pas l'absence de la fonctionnalité.
% Colonne « Assistant sur les données », trois libellés différents pour trois situations
% différentes : Onfleet a bien un assistant, mais documenté comme une aide à l'usage du produit,
% d'où « Support produit » ; Bringg n'en documente aucun, et « Non documenté » dit exactement
% cela sans affirmer qu'il n'en existe pas ; un module ERP natif, lui, n'en a pas par
% construction, d'où « Non ».
% « Sans objet » plutôt que « Non » pour les versions ERP chez Onfleet et Bringg : la question ne se
% pose pas à une plateforme qui reste sur une API générique, elle ne surgit qu'en descendant au
% modèle de données. Les marquer « Non » laisserait croire à une faiblesse là où il n'y a qu'un
% choix d'architecture différent.
\begin{tabularx}{\textwidth}{L{1.9cm}C{1.3cm}C{1.5cm}C{1.8cm}C{1.3cm}C{1.9cm}>{\raggedright\arraybackslash}X}
\hline
\rowcolor{isimsblue!10}
\textbf{Solution} & \textbf{Multi-ERP} & \textbf{Versions ERP} & \textbf{Remise de caisse} &
\textbf{Hors ligne} & \textbf{Assistant sur les données} & \textbf{Limite principale} \\
\hline
Onfleet & Partiel & Sans objet & Non & Option & Support produit & Intégration ERP superficielle \\
\hline
Bringg & Partiel & Sans objet & Partiel & Oui & Non documenté & Mise en œuvre lourde, pensée par ERP \\
\hline
Module ERP natif & Non & Non & Partiel & Non & Non & Enferme l'entreprise dans un seul ERP \\
\hline
\textbf{Solution proposée} & \textbf{Oui} & \textbf{Oui} & \textbf{Oui} & \textbf{Oui} &
\textbf{Oui} & --- \\
\hline
\end{tabularx}
\end{table}

La colonne « versions ERP » mérite un mot. Un progiciel évolue, et ses versions majeures modifient
la structure des données qu'il expose : une intégration construite sur une version peut cesser de
fonctionner chez un client qui met la sienne à jour. Les plateformes étudiées échappent à cette
difficulté en restant à la surface, à travers une API générique, d'où la mention « sans objet » ;
un module natif la subit au contraire de plein fouet, puisqu'il doit être porté à chaque version.
S'intégrer en profondeur tout en résistant aux montées de version est donc une exigence à part
entière, dont le chapitre~\ref{ch:erp} détaille le traitement.

Ce constat oriente la conception du système autour de quatre partis pris : un adaptateur ERP commun
plutôt qu'un développement par ERP, une intégration qui s'adapte à la version rencontrée chez le
client plutôt que d'en supposer une seule, une isolation des données par entreprise cliente plutôt
qu'un déploiement dédié à chacune, et une chaîne d'encaissement traitée comme un circuit de
responsabilité, où le livreur déclare la somme remise, un tiers la compte et un responsable arbitre
les écarts, plutôt que comme un simple montant saisi à la livraison.

\section{Méthodologie de développement}

Le projet a été mené par un développeur unique, sur un périmètre fonctionnel qui s'est précisé au
fil de l'avancement, notamment sur l'étendue exacte de l'intégration multi-ERP. Le choix d'une
méthode s'est donc fait sur ce que chacune exige de celui qui l'applique, plutôt que sur ses mérites
en général. Le tableau~\ref{tab:comparatif-methodo} confronte les trois approches envisagées aux
quatre contraintes réelles de ce projet.

\begin{table}[H]
\centering
\caption{Ce que chaque méthode suppose, au regard des contraintes de ce projet}
\label{tab:comparatif-methodo}
\small
\begin{tabular}{L{3.3cm}L{3.5cm}L{3.5cm}L{3.7cm}}
\hline
\rowcolor{isimsblue!10}
\textbf{Contrainte du projet} & \textbf{Cycle en V} & \textbf{RUP} & \textbf{Scrum} \\
\hline
Le périmètre s'est précisé en cours de route & Exige un cahier des charges arrêté avant de commencer
& Révision possible, au prix d'un changement de phase & Le reste à faire est réexaminé à chaque
cycle \\
\hline
Le projet est mené par une seule personne & Suppose des intervenants distincts par phase & Prévoit de
nombreux artefacts et disciplines & Trois rôles, cumulables par une même personne \\
\hline
Un résultat montrable était attendu tôt & Rien d'utilisable avant la fin & Prototype à la fin de
chaque phase & Un incrément fonctionnel par cycle \\
\hline
L'encadrement demandait un point régulier & Jalons espacés & Revues de phase & Revue et
rétrospective toutes les trois à cinq semaines \\
\hline
\end{tabular}
\end{table}

Le choix s'est porté sur la méthode agile \textbf{Scrum}, pour ses cycles courts et sa capacité à
réévaluer les priorités à chaque sprint plutôt qu'une seule fois en début de projet.

\subsection{Présentation de Scrum}

Scrum~\cite{schwaber20} découpe le projet en cycles courts de durée fixe, appelés \emph{sprints}, chacun devant
produire quelque chose d'utilisable plutôt qu'une étape intermédiaire. Ce qui reste à faire est tenu
dans une liste unique et ordonnée, le \emph{Product Backlog}, dont on détache au début de chaque
cycle ce que ce cycle peut absorber. Ce qui a rendu la méthode utile ici tient surtout à la fin de
cycle~: la revue oblige à montrer ce qui tourne réellement, et la rétrospective à corriger la façon
de travailler avant d'engager la suite. Le périmètre de ce projet s'étant précisé en cours de route,
cette réévaluation régulière valait mieux qu'un plan arrêté au départ.

\subsection{Organisation du travail et outils}

Le projet ayant été mené par un seul développeur, les trois rôles Scrum, Product Owner, Scrum
Master et équipe de développement, ont été assumés par la même personne. Le suivi du Product
Backlog et des Sprint Backlogs a été réalisé avec \textbf{GitLab Issues et Boards}, qui a servi de
tableau kanban pour chaque sprint.

Le développement a été découpé en un sprint d'initialisation (sprint 0) suivi de cinq sprints
fonctionnels, chacun clôturé par une revue et une rétrospective avant l'engagement du suivant. Le
chapitre suivant détaille le sprint 0.

\section{Conclusion}

Ce chapitre a présenté le cadre général du projet : l'entreprise d'accueil ASM, la problématique de
l'intégration ERP multiple sans dépendance ni interruption réseau, les objectifs et le périmètre
retenus, ainsi que les limites des solutions existantes qui justifient cette approche. La
méthodologie Scrum a été retenue pour sa capacité à absorber un périmètre encore mouvant. Le
chapitre suivant présente le sprint 0 : l'analyse des besoins, la planification des sprints et
l'étude technique préalable.
